\documentclass{book} % MUST be first
\usepackage{amsthm}
\newtheorem{definition}{Definition}
\newtheorem{axiom}{Axiom}% class and options
\usepackage{fontspec}
\setmainfont{Libertinus Serif}
\usepackage{unicode-math}
\setmathfont{Libertinus Math}   % load useful packages
\begin{document}
\begin{center}
\Large{High School Mathematics: An Intuitive Introduction}
\newline
\newline
\normalsize{By LinuxChad628}
\newpage
To the GOAT Mr. P (You know who you are)
\newpage
\end{center}
\chapter*{Introduction}
The goal of this book is to deliver an intuitive introduction to high-school level mathematics.
My goal here is not to replace high school, but rather to provide you with the necessary mathematical background to do well in your classes.
As for how you should read this book, that depends on you comfort level when it comes to mathematics. If you are new to this type of math, you should read slowly,
and make sure to do all of the exercises. On the other hand, if you are comfortable with math, you can go faster and only do the exercises that you think you need extra practice with.
However, no matter your level of expertise, I would recommend reading the entire thing and not skipping any of the book.
This is because the goal of this book is not only to teach you math, but also to introduce a new way of looking at this subject, one that focuses on intuition first, rather than raw computation.
\newpage
\begin{center}Contents:\end{center}
\small{Unit 1: Algebra I}
\newline
\indent \footnotesize{Chapter 1: Expressions and Equations}
\newline
\indent \footnotesize{Chapter 2: Graphs of Lines}
\newline
\indent \footnotesize{Chapter 3: Solving Linear Equations}
\newline
\indent \footnotesize{Chapter 4: Systems of Equations}
\newline
\indent \footnotesize{Chapter 5: Graphing 1-D and 2-D Inequalities}
\newline
\indent \footnotesize{Chapter 6: Solving Inequalities}
\newline
\indent \footnotesize{Chapter 7: The Quadratic}
\newline
\indent \footnotesize{Chapter 8: Solving the Quadratic By Graphing}
\newline
\indent \footnotesize{Chapter 9: Solving the Quadratic By Factorization}
\newline
\indent \footnotesize{Chapter 10: Perfect Square Quadratics}
\newline
\indent \footnotesize{Chapter 11: Completing the Square}
\newline
\indent \footnotesize{Chapter 12: The General Quadratic and the Quadratic Formula}
\newline
\indent \footnotesize{Chapter 13: Absolute Value}
\newline
\indent \footnotesize{Chapter 14: Absolute Value Equations and Inequalities}
\newline
% Unit 2: Geometry
\small{Unit 2: Geometry}
\newline
\indent \footnotesize{Chapter 15: Preliminary Notions of Geometry}
\newline
\indent \footnotesize{Chapter 16: Proof}
\newline
\indent \footnotesize{Chapter 17: Line Segments and Angles}
\newline
\indent \footnotesize{Chapter 18: Parallel Lines}
\newline
\indent \footnotesize{Chapter 19: Triangles}
\newline
\indent \footnotesize{Chapter 20: Polygons}
\newline
\indent \footnotesize{Chapter 21: Transformations}
\newline
\indent \footnotesize{Chapter 22: Extending the Real Numbers}
\newline
\indent \footnotesize{Chapter 23: Transformations in the Complex Plane}
\newline
\indent \footnotesize{Chapter 24: Trigonometry}
\newline
\indent \footnotesize{Chapter 25: Circles}
\newline
% UNIT 3: Precalculus
\small{Unit 3: Precalculus}
\newline
\indent \footnotesize{Chapter 26: Exploring Functions}
% IDK what to put here
\chapter{Expressions and Equations}

All of algebra is built on the idea of \textit{expressions} and \textit{equations}.
Expressions and equations allow us to express ideas compactly using symbols, and they will continue to be critical for the rest of your mathematical career.

\begin{definition}
Expression: An expression is a group of symbols and operators that represent a value. An expression never includes an equals sign.
\end{definition}

\begin{definition}
Equation: An equation consists of two expressions related by equality, in the form $a=b$ where $a$ and $b$ are expressions.
\end{definition}

\noindent Here are a few examples of expressions:
\[1 + 1\]
\[x + 3\]
\[3y - 17\]
\[\tau - 1\]

\noindent Here are some equations:
\[A=\frac{1}{2}\tau r^{2}\]
\[x+3=4\]
\[x=1\]
\[1+1=2\]
\[\tau=6.28318530717\dots\]

Expressions and equations can contain constants, variables, or a mix of the two. 
A variable is a quantity, often represented by a letter, whose value may change.
A constant is a fixed value.
Examples of constants include $1$, $6$, $\tau$, and $12.463$.
\begin{definition}
    \textit{Numerical Expression}: A numerical expression contains only constants
\end{definition}
\begin{definition}
    \textit{Algebraic Expression}: An expression containing at least one variable.
\end{definition}
When one is given an expression, their first goal should be to simplify the expression.
An expression is considered simplified when all possible operations have been performed.
For example, if one is given the expression $1+1$, they will usually simplify it to $2$.
Algebraic expressions may be simplified as well.
In this case, one must apply a technique called "combining like terms".
Often variables will have numbers in front of them, these numbers are called coefficients.
An example of a coefficient is the "3.5" in $3.5x$.
This expression is equal to $3.5 * x$.
When you combine like terms, you first group together all the variables of the same type.
Then, you add up their coefficients, and multiply the variable by that coefficient.
This new expression can replace all instances of that variable in the original expression.
\newline \newline Example:
\[x + 3y + 2x - 0.5y\]
\[= 2x + x + 3y - 0.5y\]
\[= 3x + 2.5y\]
Note that $x$ is the same as $1x$.
\newline \newline
\indent As you saw in definition 2, an equation consists of two expressions separated by an equality.
Therefore, if one is given an equation, they may simplify each side, and those sides will remain equal.
Example:
\[x + 2x + 1 = 3 - y + 4 + y\]
\[3x + 1 = 7\]
\indent When a single variable appears in an equation, it is called an unknown.
This is because that variable can take on only 1 value.
In the equation $3x+1=7$, $x=2$.
This can be verified by substituting $2$ for $x$:
\[3(2)+1=7\]
\[6+1=7\]
\[7=7\]
Which is true.
\newline \indent Because $x$ can only take on one value, it is called an "unknown" until its value is determined, then one may refer to it as a constant.
Much of algebra focuses on determining the values of unknowns, which we will soon learn.
\newline \indent In an equation like $y + 1 = 3x -7$, $x$ may take on many values (infinitely many in fact), and so may $y$.
That is why, in this context, we refer to $x$ and $y$ and \textit{variables}.
\Large{Exercises:}
1. Determine whether each of the following is an equation or an expression:



\end{document}
